\scanpage{27}
\deleted{$\Theta_\varphi$, la forme $\mathrm{Tr}\,uv$ devient
la forme sur $\Phi=\unsure{\Gamma|\mathscr T}(M^\vee\otimes M^\vee(p))$ \unsure{de} induite
par $\overset{\vee}{\varphi}\otimes\overset{\vee}{\varphi}(p)$, ou encore,}
identifiant $A$ d'une part à
$\Phi(M)=M^\vee\otimes M^\vee(p)$ \inserted{par $\Theta_\varphi$}, d'autre part via
\deleted{\unsure{$\varphi$}} $u\mapsto{}^tu$ et $\overset{\vee}{\varphi}$ à [\unsure{ou|où} via $\Phi(M)\simeq\Phi(M^{*})$
à l'aide de l'isom $\widetilde\varphi:M\longrightarrow M^*$] à
$\Phi(M^*)=M\otimes M(-p)$, la forme $\mathrm{Tr}\,uv$
sur $A\times A$ devient la forme évidente
de contraction sur $\Phi(M)\times\Phi(M^*)$.

Noter d'autre part que \deleted{pour} \inserted{\unsure{à}} tt ``cône
\deleted{positif sur} admissible'' $\Phi^+(M)$ de $\Phi(M)$,
est canoniquement associé un ``cône admissible''
$\Phi^+(M^*)$ de $\Phi(M^*)$, en prenant les
$\overset{\vee}{\varphi}^*$ ($\varphi\in\Phi^{\unsure{+|*}}(M^*)^\circ$). Ceci posé, on voit
que \uline{$\Phi^+(M)$ et $\Phi^+(M^*)$} sont
polaires l'un de l'autre. Conséquence :
% archive sheet number: 26
